\chapter{Итоговая электромагнитная проверка барионной ветви}
\label{ch:v8-baryon-electromagnetic-closure-redteam-gate}

Сведём зарядовое тождество, пространственное усреднение, магнитный контакт
и перестановочные ветви в одну границу утверждений.

\section{Общая условная формула}

Обозначим положительную электростатическую комбинацию
\begin{equation}
 A_{\rm el}=\mu+2\bar g>0,\qquad z=\zeta h.
 \label{eq:v8-baryon-em-closure-parameters}
\end{equation}
Для трёх согласованных перестановочных типов предыдущих глав полная
электромагнитная разность имеет вид
\begin{equation}
 3T(E_n-E_p)=
 \begin{cases}
  -A_{\rm el}-z,&\mathbf1,\\
  -A_{\rm el}+3z,&\mathbf1_{\rm sgn},\\
  -A_{\rm el}+z,&\mathbf2.
 \end{cases}
 \label{eq:v8-baryon-em-three-branch-total-splitting}
\end{equation}
Отрицательный знак в каждой ветви эквивалентен соответственно
\begin{equation}
 z>-A_{\rm el},\qquad
 z<\frac{A_{\rm el}}3,\qquad
 z<A_{\rm el}.
 \label{eq:v8-baryon-em-branch-sign-conditions}
\end{equation}
Поэтому знак отрицателен одновременно во всех трёх ветвях только в полосе
\begin{equation}
 -A_{\rm el}<z<\frac{A_{\rm el}}3.
 \label{eq:v8-baryon-em-common-negative-strip}
\end{equation}
Текущая теория не выводит ни ветвь, ни число $z$, ни такую оценку.

\section{Точные контрмодели без смены электростатического сектора}

При фиксированном $A_{\rm el}>0$ знак можно изменить в каждой ветви:
\begin{equation}
 \begin{array}{c|c|c}
 \text{ветвь}&z&3T(E_n-E_p)\\ \hline
 \mathbf1&-2A_{\rm el}&A_{\rm el}\\
 \mathbf1_{\rm sgn}&A_{\rm el}&2A_{\rm el}\\
 \mathbf2&2A_{\rm el}&A_{\rm el}
 \end{array}.
 \label{eq:v8-baryon-em-closure-sign-countermodels}
\end{equation}
Это не утверждение о реализуемости данных значений конкретной квантовой
электродинамикой. Это сертификат логической недостаточности имеющихся
аксиом: магнитный коэффициент и пространственный контакт ими не ограничены.

\section{Что закрыто и что остаётся внешним}

Строго закрыты:
\begin{enumerate}
 \item операторное тождество $A+C=Q_{\rm tot}^2$;
 \item зарядовый рисунок $(4,1,0,1)$;
 \item условный электростатический знак после перестановочного усреднения;
 \item три перестановочно-допустимые магнитные ветви;
 \item масштабные орбиты кулоновского и контактного матричных элементов.
\end{enumerate}

Не выведены: координатный гамильтониан, радиальная волновая функция,
физический спиновый носитель исходной модели, магнитные моменты, кварковые
массы, коэффициент $z$, выбор перестановочной ветви и абсолютная единица
энергии. Поэтому итоговый физический знак и величина разности масс не
являются следствием Тома VIII.

\begin{nogocriterion}[Закрытие электромагнитной ветви]
Формулу полного заряда нельзя повышать до параметр-свободного
электромагнитного предсказания. Для продолжения требуется новый
трёхкоординатный родитель, который одновременно выбирает радиальный
масштаб, перестановочную ветвь и магнитный контакт.
\end{nogocriterion}

\begin{protocol}
\begin{enumerate}
 \item зарядовая алгебра: \textbf{закрыта точно};
 \item условный электростатический знак: \textbf{закрыт};
 \item магнитные перестановочные ветви: \textbf{классифицированы};
 \item общая полоса отрицательного знака:
       \textbf{$-A_{\rm el}<z<A_{\rm el}/3$};
 \item принадлежность физической модели этой полосе: \textbf{не выведена};
 \item полный электромагнитный знак: \textbf{не фиксирован};
 \item абсолютная величина: \textbf{не фиксирована};
 \item электромагнитная ветвь Тома VIII: \textbf{закрыта отрицательно};
 \item условие возобновления: \textbf{координатно-спиновый родитель}.
\end{enumerate}
\end{protocol}

Машинный сертификат сохранён в
\path{s2t_v8_baryon_electromagnetic_closure_redteam_gate_results.json}.